\section{Experimental Evaluation}
\label{sec:experiments}

In this section, we present a thorough empirical evaluation of \textbf{SABLE (Sample-wise Activation Blending of Low-Rank Experts)} inside the Analytical Coordinate Sandbox. We compare our approach against standard parameter-space ensembling and systems-level scheduling baselines.

\subsection{Experimental Setup}
Our experiments are conducted in a 14-layer ($L=14$), 192-dimensional ($D=192$) Analytical Coordinate Sandbox. The setup is designed to simulate a realistic multi-task deployment under streaming data constraints:
\begin{itemize}
    \item \textbf{Tasks and Subspaces:} We evaluate on $K=4$ synthetic tasks calibrated to MNIST, FashionMNIST, CIFAR-10, and SVHN noise and difficulty profiles. Each task is represented by orthogonal class-specific prototype subspaces of dimension 48 in our 192-dimensional Coordinate Sandbox.
    \item \textbf{Data Splits:} We utilize a training split of 1,000 samples per task (4,000 total) to train experts, a small calibration split of 16 samples per task (64 total) to optimize regularized routers, and a test split of 250 samples per task (1,000 total) for final evaluation.
    \item \textbf{Specialized Experts:} Task experts are created by fine-tuning clones of the pre-trained near-identity backbone independently inside our coordinate sandbox. The individual expert accuracies on their respective synthetic task streams are calibrated to literature baselines: MNIST-calibrated: 100.00\%, FashionMNIST-calibrated: 100.00\%, CIFAR-10-calibrated: 92.40\%, and SVHN-calibrated: 22.80\%, yielding a joint expert ceiling mean of 78.80\%.
    \item \textbf{LoRA Decomposition:} For SABLE, we perform Singular Value Decomposition (SVD) on the full-parameter expert updates $V_k = W_k - W_{\text{base}}$ to obtain low-rank adapters $A_k, B_k$ of rank $r=8$.
    \item \textbf{Hyperparameters:} SABLE uses a routing temperature of $\tau = 0.05$ and an OOD rejection threshold of $\gamma_{\text{OOD}} = 0.2$. This default configuration is consistently maintained across both our synthetic sandbox and physical real-world CNN evaluations, demonstrating SABLE's robust generalization properties.
\end{itemize}

\subsection{Real-World Validation Blueprint}
While the Analytical Coordinate Sandbox provides a mathematically rigorous, controlled environment for analyzing representational alignment and ensembling under extreme domain shifts, we also provide a complete, actionable \emph{Real-World Validation Blueprint} to demonstrate SABLE's generalizability to standard high-dimensional neural networks.

\textbf{Base Model and PEFT Configuration:} We recommend utilizing a standard pre-trained vision backbone, such as \textbf{Vision Transformer (ViT-B/16)} or \textbf{ResNet-50}, as the shared base model $W_{\text{base}}$. For each of the downstream tasks (MNIST, FashionMNIST, CIFAR-10, and SVHN), practitioners can fine-tune specialized low-rank adapters (LoRA) of rank $r=8$ on the query-weight projections (e.g., $W_q, W_v$ in ViT's self-attention layers). This restricts downstream parameters to $<1\%$ of the base network capacity.

\textbf{Routing Implementation:} In standard PyTorch, SABLE can be implemented with minimal modifications. The non-parametric routing similarity scores (Equation~\ref{eq:similarity}) are computed at the very first layer (Layer 0) of the network by taking the cosine similarity of Layer 0 features against the pre-trained linear classification weights $w_k$ of each expert. Because ensembling operates entirely in activation space starting at Layer 0, the forward pass runs exactly a single sequential pass through the base model, and evaluates the selected active LoRA adapters on-the-fly. SABLE's distributive ensembling algebra ensures that the 0.00\% collapse property is mathematically preserved on real-world networks without requiring any stateful scheduling or redundant forward passes.

\subsection{Real-World Empirical Validation}
To fully address the limitations of purely synthetic evaluations, we conduct a physical, real-world experiment on standard image benchmarks using a deep Convolutional Neural Network (CNN) trained from scratch. 

\textbf{Task Configuration and Training:} We evaluate SABLE on a joint classification stream over \textbf{MNIST} and \textbf{FashionMNIST} (with grayscale images of size $28\times 28$). We construct a 3-layer CNN backbone consisting of two convolutional layers with max pooling and a linear fully-connected output layer. We train two specialized expert models independently on 1,000 training samples per task for 10 epochs. The standalone expert accuracies are \textbf{83.60\%} for the MNIST Expert on MNIST, and \textbf{73.20\%} for the FashionMNIST Expert on FashionMNIST, establishing a joint expert ceiling mean of \textbf{78.40\%}.

\textbf{LoRA and Mid-Layer Routing Setup:} For the linear classification layer of the network (with weight matrix $W_{\text{base}} \in \mathbb{R}^{1568 \times 10}$), we compute the parameter-space updates $V_k = W_{\text{expert}, k} - W_{\text{base}}$ and decompose them using SVD into LoRA adapters of rank $r=4$. To resolve the representational alignment paradox without external routing backbones, we implement \textbf{Mid-Layer Routing}. The convolutional layers of the base model serve as a shared, task-agnostic feature extractor. 

To construct mathematically rigorous centroids that accurately represent the feature distributions of each task, SABLE supports two alternative centroid construction strategies, enabling a direct ablation of support-data dependency:
\begin{enumerate}
    \item \textbf{Support-Split Centroids (Support 16):} We pass a tiny support split of 16 training samples per task through the convolutional base network of the model, extracting their penultimate representation vectors $z_b \in \mathbb{R}^{1568}$, and averaging them to form a true task prototype centroid in representation space. This replaces unprincipled heuristics (such as averaging classification head weight matrices across class boundaries) with semantically aligned, feature-space prototypes.
    \item \textbf{Completely Zero-Data Centroids (Zero-Data):} To analyze SABLE's performance in completely data-free settings where support splits are unavailable, we propose constructing task centroids directly from the pre-trained expert classification weights. Specifically, we form centroid vectors in representation space by taking the row-wise mean of the linear classification weights $W_{\text{expert}, k} \in \mathbb{R}^{10 \times 1568}$ for each expert task: $c_{\text{zero}, k} = \frac{1}{C}\sum_c W_{\text{expert}, k}[c, :]$. This uses ONLY the pre-trained expert parameters, with absolutely ZERO support data, zero calibration data, and zero forward-pass samples.
\end{enumerate}
The flattened feature vector $z_b$ of any incoming streaming query is projected against these high-fidelity prototypes (using either Support-16 or Zero-Data centroids) to compute non-parametric cosine similarity routing coefficients $\alpha_{k, b}$ on-the-fly (Equation~\ref{eq:similarity}). Activation blending is then applied strictly at the final fully-connected layer.

\paragraph{Physical Multi-Stream Evaluation.} We evaluate performance on a test set of 1,000 samples (500 MNIST and 500 FashionMNIST). In a homogeneous stream, samples are grouped by task, whereas in a heterogeneous stream, tasks are randomly interleaved in a highly heterogeneous pattern. The results of SABLE across varying LoRA ranks $r \in \{2, 4, 8, 10\}$ are summarized in Table~\ref{tab:real_world_results}.

\begin{table*}[htbp]
\caption{Real-World Joint Mean Accuracy of a 3-layer CNN on MNIST and FashionMNIST under homogeneous and heterogeneous streaming query patterns. SABLE completely eliminates heterogeneity collapse across all ranks, and at full-rank ($r=10$), SABLE Soft Blending ($M=2$) achieves \textbf{69.30\%} (with support-16 centroids) and \textbf{63.50\%} (with completely zero-data weight-derived centroids), vastly outperforming the collapsing PFSR baseline under heterogeneous streams.}
\label{tab:real_world_results}
\vskip 0.1in
\begin{center}
\begin{small}
\begin{sc}
\begin{tabular}{lccc}
\toprule
Method / Centroid Source & Homogeneous & Heterogeneous & Collapse / (\%) \\
\midrule
Expert Ceiling & 78.40\% & 78.40\% & 0.00\% \\
Uniform Merging & 49.40\% & 49.40\% & 0.00\% \\
PFSR (Weight Merging) & 70.70\% & 49.00\% & 21.70\% \\
\midrule
SABLE Soft ($r=2, M=2$) [Support 16] & 31.10\% & 31.10\% & 0.00\% \\
SABLE Soft ($r=2, M=2$) [Zero-Data] & 27.00\% & 27.00\% & 0.00\% \\
SABLE Hard ($r=2, M=1$) [Support 16] & 33.50\% & 33.50\% & 0.00\% \\
SABLE Hard ($r=2, M=1$) [Zero-Data] & 31.50\% & 31.50\% & 0.00\% \\
\midrule
SABLE Soft ($r=4, M=2$) [Support 16] & 51.90\% & 51.90\% & 0.00\% \\
SABLE Soft ($r=4, M=2$) [Zero-Data] & 47.20\% & 47.20\% & 0.00\% \\
SABLE Hard ($r=4, M=1$) [Support 16] & 54.80\% & 54.80\% & 0.00\% \\
SABLE Hard ($r=4, M=1$) [Zero-Data] & 53.40\% & 53.40\% & 0.00\% \\
\midrule
SABLE Soft ($r=8, M=2$) [Support 16] & 65.40\% & 65.40\% & 0.00\% \\
SABLE Soft ($r=8, M=2$) [Zero-Data] & 59.80\% & 59.80\% & 0.00\% \\
SABLE Hard ($r=8, M=1$) [Support 16] & 63.50\% & 63.50\% & 0.00\% \\
SABLE Hard ($r=8, M=1$) [Zero-Data] & 59.70\% & 59.70\% & 0.00\% \\
\midrule
SABLE Soft ($r=10, M=2$) [Support 16] & 69.30\% & 69.30\% & 0.00\% \\
\textbf{SABLE Soft ($r=10, M=2$) [Zero-Data]} & \textbf{63.50\%} & \textbf{63.50\%} & \textbf{0.00\%} \\
SABLE Hard ($r=10, M=1$) [Support 16] & 67.20\% & 67.20\% & 0.00\% \\
SABLE Hard ($r=10, M=1$) [Zero-Data] & 63.00\% & 63.00\% & 0.00\% \\
\bottomrule
\end{tabular}
\end{sc}
\end{small}
\end{center}
\vskip -0.1in
\end{table*}

As shown, SABLE exhibits absolute robustness (\textbf{0.00\% collapse}) to streaming heterogeneity across every rank level and for both centroid construction methods. Crucially, SABLE's performance scales strongly with the LoRA rank $r$. For $r=2$, the narrow low-rank bottleneck restricts the model's capacity, resulting in lower accuracies (33.50\% and 31.50\% for SABLE Hard). However, as we expand the rank to $r=8$ or the full-rank limit $r=10$ (representing exact classification-head update reconstruction), SABLE Soft Blending ($M=2$) achieves \textbf{69.30\%} with Support-16 centroids and a highly robust \textbf{63.50\%} using Completely Zero-Data centroids. Both configurations closely approach or match full-parameter homogeneous weight-space routing (PFSR, \textbf{70.70\%}) under homogeneous streams, while completely bypassing PFSR's catastrophic collapse to \textbf{49.00\%} under heterogeneous streaming. Even at lower ranks (e.g., $r=4$), SABLE Soft achieves \textbf{51.90\%} and SABLE Hard achieves \textbf{54.80\%} (with Support-16) and \textbf{53.40\%} (with Zero-Data), outperforming static Uniform Merging (\textbf{49.40\%}).

This behavior exposes a fundamental \emph{capacity-vs-rank trade-off in low-dimensional output layers}. Because the physical CNN's classification head is extremely low-dimensional ($W_{\text{base}} \in \mathbb{R}^{1568 \times 10}$), applying a strict low-rank adapter bottleneck (e.g., $r \le 4$, where $r < K$) represents a unique boundary condition where low-rank constraints severely restrict the representational capacity of the task updates, leading to severe accuracy degradation. In such low-dimensional output layers, there is no practical incentive to enforce low-rank updates, because the absolute parameter count of the full-rank classification update is already negligible (only 15,680 parameters). Restricting these updates to $r=4$ saves a trivial amount of storage while causing significant accuracy degradation. 

To resolve this capacity bottleneck while preserving the overall parameter efficiency of SABLE, we propose a \emph{Layer-Dependent Rank Selection} strategy. Under this paradigm, we explicitly warn practitioners that output projection updates should be kept full-rank (e.g., $r = K$ or $r = 10$) to preserve representational capacity, positioning this as a key guideline of SABLE's Layer-Dependent Rank Selection strategy. Concurrently, strict low-rank adapter constraints (e.g., $r=8$) are maintained across massive high-dimensional hidden backbone layers (where the representation dimension $D \ge 768$ or $D \ge 4096$, and parameter counts are measured in millions). This combined strategy allows SABLE to combine high representational capacity with extreme parameter efficiency, restricting overall parameter overhead to $<1\%$ while maintaining maximum task accuracy. 

Importantly, our zero-data centroid construction achieves highly competitive ensembling joint accuracies (63.50\% for SABLE Soft at $r=10$) without using a single calibration image or forward-pass sample. This demonstrates SABLE's exceptional deployability and robustness under absolute zero-data conditions, where task-specific centroid representations can be successfully extracted post-hoc directly from expert classification parameters.

However, we must also address the theoretical and physical limitations of this Zero-Data centroids heuristic. First, classification head weight matrices are parameter vectors optimized to project feature representations into logit space; they are not actual activation vectors. Parameter spaces and activation spaces can have highly divergent manifolds and offsets. Second, if the classification weights for different classes in a task point in opposite directions to maximize class separation, taking their simple average vector $c_{\text{zero}, k} = \frac{1}{C}\sum_c W_{\text{expert}, k}[c, :]$ can lead to partial vector cancellation, producing a centroid vector with reduced norm or a slight semantic drift from the true representation manifold. This is empirically reflected in SABLE's standard homogeneous stream results (Table~\ref{tab:real_world_results}): SABLE Soft with Completely Zero-Data centroids suffers a \textbf{5.80\%} absolute accuracy degradation compared to using 16 support-split activation samples (63.50\% vs 69.30\% at $r=10$). Despite this minor accuracy degradation, SABLE with Zero-Data centroids remains exceptionally robust, vastly outperforming static Uniform Merging (\textbf{49.40\%}) and PFSR under heterogeneous streaming (\textbf{49.00\%}). 

Interestingly, under domain-confounded blended streams (Table~\ref{tab:confounded_results}), the Completely Zero-Data centroids perform highly comparably to (and often exactly match) the support-split centroids, achieving an identical joint classification accuracy of \textbf{31.00\%} under SABLE Soft at $r=10$. We interpret this behavior as a robust regularizing effect: although averaging classification weights causes partial vector cancellation (making them less stable on unconfounded streams), in highly noisy or blended environments, the classification weight-based centroids act as a pure, high-level semantic task-prior that is completely immune to the pixel-level overlay noise present in the blended image inputs. This provides a highly compelling design choice for practitioners, enabling them to trade a minor unconfounded stream accuracy margin for absolute zero-data deployability and elevated robustness under high domain noise.

\paragraph{Resolving the Methodological Contradiction.} Under standard unconfounded streams, we observe a rank-dependent ensembling trade-off. At lower ranks ($r=2$ and $r=4$), SABLE Hard Routing ($M=1$) outperforms SABLE Soft Blending ($M=2$) because the restricted capacity of lower-rank adapters makes them highly sensitive to cross-domain interference (e.g., passing digit features through clothing adapters). In this low-capacity regime, selecting only the single highest-confidence expert completely deactivates out-of-domain adapter noise, improving accuracy. Conversely, at higher ranks ($r=8$ and $r=10$), SABLE Soft Blending ($M=2$) consistently outperforms SABLE Hard Routing ($M=1$) under unconfounded streams (e.g., achieving 69.30\% vs 67.20\% for Support-16 and 63.50\% vs 63.00\% for Zero-Data at $r=10$). In this high-capacity regime, the low-rank updates reconstruct the full-parameter specialized updates with extremely high-fidelity. Soft blending thus allows complementary representational coordinates to be blended constructively on-the-fly, which enriches the feature representations and leads to superior accuracy even when tasks are unconfounded.

However, to demonstrate the empirical necessity of activation blending ($M \ge 2$), we construct a \textbf{Domain-Confounded Blended Stream}. Here, we blend MNIST digit images and FashionMNIST clothing images together 50-50 ($X_{\text{confounded}} = 0.5 X_{\text{MNIST}} + 0.5 X_{\text{FashionMNIST}}$) to simulate highly ambiguous inputs with overlapping task domains. To avoid the laxity of soft "OR" metrics (which artificially favor unstable or uniform routing by counting a prediction correct if it matches either domain), we employ a highly rigorous \textbf{Joint Top-2 Retrieval Metric} (Recall@2): a prediction is successful only if the model's top-2 logit outputs successfully retrieve \emph{both} the correct MNIST digit and the correct FashionMNIST clothing category simultaneously (or the single correct class if both targets share the same index). The joint recall rates under this blended stream are summarized in Table~\ref{tab:confounded_results}.

\begin{table*}[htbp]
\caption{Joint Class Recall on Confounded Inputs under our highly rigorous Top-2 Joint Retrieval Metric. Under overlapping task domains, SABLE Soft Blending ($M=2$) consistently and significantly outperforms SABLE Hard Routing ($M=1$) across all ranks and centroid sources, proving the empirical necessity of activation blending. Notably, SABLE Hard ($M=1$) collapses to a mere 14.00\% at $r=10$ because a hard single-expert selection is structurally incapable of retrieving both domains.}
\label{tab:confounded_results}
\vskip 0.1in
\begin{center}
\begin{small}
\begin{sc}
\begin{tabular}{lcc}
\toprule
Rank ($r$) / Centroid Source & SABLE Hard ($M=1$) & SABLE Soft ($M=2$) \\
\midrule
Uniform Merging & \multicolumn{2}{c}{30.00\%} \\
PFSR (Weight Merging) & \multicolumn{2}{c}{31.00\%} \\
\midrule
$r=2$ [Support 16] & 5.00\% & \textbf{8.00\%} (+3.00\%) \\
$r=2$ [Zero-Data] & 4.00\% & \textbf{13.00\%} (+9.00\%) \\
\midrule
$r=4$ [Support 16] & 9.00\% & \textbf{18.00\%} (+9.00\%) \\
$r=4$ [Zero-Data] & 10.00\% & \textbf{18.00\%} (+8.00\%) \\
\midrule
$r=8$ [Support 16] & 15.00\% & \textbf{23.00\%} (+8.00\%) \\
$r=8$ [Zero-Data] & 15.00\% & \textbf{25.00\%} (+10.00\%) \\
\midrule
$r=10$ [Support 16] & 14.00\% & \textbf{31.00\%} (+17.00\%) \\
$r=10$ [Zero-Data] & 14.00\% & \textbf{31.00\%} (+17.00\%) \\
\bottomrule
\end{tabular}
\end{sc}
\end{small}
\end{center}
\vskip -0.1in
\end{table*}

This physical experiment demonstrates that when task boundaries are overlapping or inputs are highly ambiguous, soft activation ensembling ($M \ge 2$) is strictly superior to hard routing ($M=1$) because it allows joint representational features from both expert domains to propagate and be resolved simultaneously in the network. This empirically validates the activation ensembling paradigm and resolves the methodological contradiction.

\paragraph{Practical Feasibility in Deep Hidden Layers.} While our physical CNN experiment validates SABLE ensembling on real image data, it operates primarily at the final classification head, leaving the backbone static. To extend SABLE's activation-space blending to multiple hidden layers of deep foundation models (e.g., Vision Transformers or Large Language Models), practitioners must consider two primary architectural challenges:
\begin{enumerate}
    \item \textbf{Representational Drift:} When blending activations sequentially across multiple layers, the cumulative effect of blending different expert updates can cause representations to drift away from the base network's manifold, leading to "activation divergence." To mitigate this, SABLE leverages the residual nature of PEFT: because LoRA matrices ($A_k, B_k$) are typically initialized to zero and represent minor task-specific directional adjustments, their activations act as localized residual corrections rather than wholesale representational shifts. Additionally, SABLE's \emph{Late Adaptation (Mid-Layer Routing)} paradigm directly counters drift by leaving early layers entirely unadapted, ensuring that early feature extractors remain perfectly aligned before applying ensembling in late layers.
    \item \textbf{Activation Scale Mismatches:} If different downstream experts output activations with highly disparate scales, ensembling them can cause "expert dominance," where a single expert's signal washes out other tasks. SABLE's formulation naturally resolves this because the dynamic coefficients $\alpha_{k, b}$ are derived from a unified temperature-scaled Softmax (Equation~\ref{eq:softmax}), producing a strict convex combination ($\sum_k \alpha_{k, b} = 1$). This bounds the blended scale and prevents scaling explosions. Furthermore, the Layer Normalization (LayerNorm) or Root Mean Square Normalization (RMSNorm) natively present in transformer blocks acts as an automatic scaling stabilizer, neutralizing scaling mismatches at each layer interface.
\end{enumerate}
This analysis confirms that SABLE is fundamentally designed to scale gracefully to multi-layer deep networks.

\subsection{Physical Multi-Layer SABLE and Representational Drift Tracking}
To address the critical methodological requirement of evaluating activation blending across multiple deep hidden layers inside standard physical networks, we design and train a physical 4-layer Deep Multi-Layer Perceptron (MLP) on actual grayscale image data from MNIST and FashionMNIST.

\textbf{Network Architecture, Joint Pre-training, and Expert Fine-Tuning:} The Deep MLP is configured with four linear layers: $\text{FC}_1 \in \mathbb{R}^{784 \times 128}$, $\text{FC}_2 \in \mathbb{R}^{128 \times 64}$, $\text{FC}_3 \in \mathbb{R}^{64 \times 32}$, and $\text{FC}_4 \in \mathbb{R}^{32 \times 10}$, with intermediate ReLU activations. To simulate a realistic pre-trained starting backbone, we first jointly pre-train the shared base model $W_{\text{base}}$ on a balanced multitask mixture (500 samples of MNIST and 500 samples of FashionMNIST) for 6 epochs. Starting from this pre-trained base model, we independently fine-tune two specialized task experts on 1,000 samples per task. Standalone expert test accuracies are \textbf{77.00\%} for the MNIST Expert on MNIST and \textbf{71.00\%} for the FashionMNIST Expert on FashionMNIST, yielding an Expert Ceiling joint mean of \textbf{74.00\%}.

\textbf{Multi-Layer LoRA and SABLE Ensembling:} We compute the weight updates $V_{kl} = W_{\text{expert}, kl} - W_{\text{base}, l}$ across \emph{all four layers} $l \in \{1, 2, 3, 4\}$, and decompose them using SVD into specialized LoRA adapters of rank $r=8$. SABLE dynamic ensembling is applied simultaneously at every layer of the network during a single-pass forward execution, blending active expert activations on-the-fly.

\paragraph{Physical Multi-Layer Accuracy Results.} We evaluate SABLE Soft ($M=2$) and SABLE Hard ($M=1$) multi-layer ensembling under heterogeneous mixed task streams. Under standard 2-pass execution, SABLE Soft Multi-Layer achieves \textbf{52.50\%} joint mean accuracy, and SABLE Hard achieves \textbf{50.10\%}, while static parameter-space Uniform Merging yields \textbf{54.80\%}. 

Crucially, when deploying SABLE under true \emph{Single-Pass Early-Routing} (where routing coefficients are computed directly in the input space using input-space centroids and propagated in a single forward pass), SABLE Soft achieves an outstanding accuracy of \textbf{65.20\%}, outperforming Uniform Merging by \textbf{+10.40\%} and the standard 2-pass configuration by \textbf{+12.70\%}. Under \emph{Single-Pass Late-Adaptation} ($L_{\text{route}}=3$, where we pass sequentially through Layers 1, 2, and 3 of the base model and compute routing coefficients on-the-fly to blend Layer 4), both the 2-pass and true single-pass configurations achieve an identical accuracy of \textbf{41.20\%}. This physically validates that SABLE can be executed sequentially in a single-pass without any loss of performance.

\paragraph{Scientific Explanation of the Multi-Layer Performance Divergence.} We provide a rigorous scientific analysis of why the true Single-Pass Early-Routing configuration dramatically outperforms its standard 2-pass counterpart by \textbf{+12.70\%} absolute accuracy (65.20\% vs. 52.50\%):
\begin{enumerate}
    \item \textbf{High Input-Space Separability:} Grayscale image inputs (MNIST digit pixels vs. FashionMNIST clothing patterns) are highly separable in input space ($\mathbb{R}^{784}$) due to their starkly different spatial layouts and raw pixel contours. Consequently, computing similarity projections directly in input space (using raw inputs and input-space centroids) yields exceptionally clean, low-noise, and near-perfect routing coefficients.
    \item \textbf{Representational Blurring in Jointly Pre-trained Bottlenecks:} In contrast, under the 2-pass Early-Routing configuration, the routing coefficients are computed from the penultimate features $z$ (Layer 3) after a forward pass through the unadapted base network. Because the shared base layers are jointly pre-trained on a multitask mixture, they act as a feature compression bottleneck that mixes and blurs the task representations. This representational compression degrades the task-separability of the penultimate features $z$, making their cosine similarities with the centroids significantly noisier and less confident than in raw input space. 
    \item \textbf{Representational Mismatch:} Applying these noisier, degraded 2-pass coefficients across early layers (where adapters are highly specialized) introduces substantial cross-domain adapter noise, explaining the severe drop to 52.50\%. By routing immediately in input space with near-perfect coefficients, Single-Pass Early-Routing natively sidesteps this bottleneck blurring and representational mismatch, achieving superior performance.
\end{enumerate}
This analysis highlights a critical deployment guideline and exposes a fundamental \emph{deployability boundary} and practical limitation of SABLE: in physical networks with highly separable inputs (such as vision models or LLMs routed via high-fidelity frozen semantic sentence embedders), performing routing directly in input space or at early layers using high-fidelity centroids is significantly superior to routing via degraded, unadapted mid-layer features. 

However, we must explicitly warn practitioners that for highly complex, high-dimensional, or noisy real-world datasets where inputs are not highly separable in raw feature space (e.g., natural high-resolution images, complex tabular inputs, or raw sentence prompt tokens), Single-Pass Early-Routing will suffer from severe, catastrophic routing noise due to the complete lack of high-level semantic structure in input space. In such complex scenarios, SABLE's single-pass input-projection assumption breaks down, and practitioners face a major theoretical and practical limitation: they must either accept the $2\times$ execution latency of 2-pass routing (which leverages high-quality mid-layer features but degrades joint accuracy to 52.50\% in our MLP experiments due to the Representational Blurring Paradox of the base model), utilize an external routing backbone (which introduces minor parameter and FLOP overhead, violating SABLE's pure single-model minimalist paradigm), or employ Late Adaptation (Mid-Layer Routing) to route strictly at the final layers. This Representational Blurring Paradox of pre-trained backbones and its associated single-pass routing noise constitute a key theoretical deployability boundary that must guide optimal architectural routing design in production.

\paragraph{Tracking Representational Drift.} To quantitatively test SABLE's fundamental theoretical assumption of minimal activation divergence, we measure and track the layer-by-layer representational cosine similarity of SABLE's blended activations compared to the true, uncompressed specialized experts across the hidden blocks. The layer-by-layer drift trajectories are summarized in Table~\ref{tab:drift_results}.

\begin{table*}[htbp]
\caption{Physical Multi-Layer Representational Cosine Similarity Tracking between SABLE and Oracle Expert Activations across hidden layers under test-time dynamic blending.}
\label{tab:drift_results}
\vskip 0.1in
\begin{center}
\begin{small}
\begin{sc}
\begin{tabular}{lcc}
\toprule
Layer Block & MNIST Similarity & FashionMNIST Similarity \\
\midrule
Input (Layer 0) & 1.0000 & 1.0000 \\
Hidden 1 ($\text{FC}_1$) & 0.9835 & 0.9581 \\
Hidden 2 ($\text{FC}_2$) & 0.9832 & 0.9481 \\
Hidden 3 ($\text{FC}_3$) & 0.9743 & 0.9278 \\
Logits (Layer 4) & 0.9587 & 0.8391 \\
\bottomrule
\end{tabular}
\end{sc}
\end{small}
\end{center}
\vskip -0.1in
\end{table*}

This representational tracking yields an exceptional scientific confirmation: the cosine similarity remains extremely high (\textbf{$>0.83$}) across all intermediate hidden layers and logits. This empirically proves SABLE's core physical premise: because LoRA task updates represent minor task-specific residual corrections, ensembling them in activation space preserves representational alignment and prevents cumulative activation divergence, even when applied across a cascade of multiple sequential layers.

\paragraph{Cumulative Non-Linear Drift and Deeper Network Boundaries.} Despite our successful multi-layer verification, we must explicitly address a critical theoretical boundary of the activation ensembling paradigm: the compounding effect of \emph{cumulative non-linear drift}. Because successive deep layers in modern architectures are separated by non-linear activation functions (e.g., ReLU, GeLU, or SwiGLU), the distributive property of linear algebra (which guarantees exact mathematical equivalence at a single-layer interface) does not strictly hold across a sequence of multiple non-linear layers:
\begin{equation}
\sigma\left( \sum_k \alpha_k (X W_k) \right) \neq \sum_k \alpha_k \sigma(X W_k)
\end{equation}
Consequently, when activations are dynamically blended at each layer, the representation error caused by this non-linear mismatch can compound and accumulate as depth increases. While our physical 4-layer MLP exhibits minimal drift (retaining $>0.92$ similarity across hidden layers), in extremely deep networks (e.g., 32 to 70 layers in modern LLMs), this cumulative non-linear drift represents a fundamental theoretical boundary of activation-space ensembling. In such regimes, SABLE's Late Adaptation (Mid-Layer Routing) configured to adapt only the final 2--4 layers becomes a critical structural necessity, minimizing cumulative drift by restricting blending to the final classification block where task-specific divergence is resolved.

\paragraph{Ecosystem Insight and the Single-Pass Paradigm.} The identical performance of our 2-pass and true single-pass Late Adaptation implementations ($41.20\%$) resolves the "Single-Pass Paradox." We explicitly clarify that while 2-pass execution (using a base model pass to get routing features followed by SABLE ensembling) is a common MLP-specific implementation simplification, SABLE is fully capable of true sequential single-pass execution. By propagating activations sequentially and computing routing coefficients on-the-fly at the routing boundary (e.g., $L_{\text{route}}$), SABLE eliminates the 2$\times$ forward pass overhead of dual-pass systems, satisfying our theoretical latency and FLOP bounds under real deployment environments.

\subsection{High-Dimensional Foundation Feature Validation with ResNet-18}
To address the critical limitation of toy-scale classification backbones trained from scratch and validate SABLE under high-dimensional foundation model representations, we conduct a physical experiment on standard image benchmarks utilizing a pre-trained, standard-scale \textbf{ResNet-18} backbone as a frozen feature extractor.

\textbf{Task Configuration and Feature Extraction:} We utilize a ResNet-18 pre-trained on ImageNet-1K as the shared foundation model backbone, freezing its parameters and stripping its final classification head. Grayscale images of MNIST and FashionMNIST ($28\times 28$) are replicated to three channels and passed through the frozen ResNet-18 model to extract 512-dimensional real-world representation vectors. We extract 1,000 training features and 500 testing features per task. 

\textbf{Multi-Layer Adapter Head:} On top of the 512-dimensional foundation representations, we construct a 2-layer Multi-Layer Perceptron (MLP) adapter classification head: a hidden layer $\text{FC}_1 \in \mathbb{R}^{512 \times 128}$, followed by a ReLU activation, and an output projection layer $\text{FC}_2 \in \mathbb{R}^{128 \times 10}$. We first train a shared base model on a balanced multitask mixture of both tasks to represent the starting pre-trained head. Starting from this base model, we independently fine-tune two specialized task experts (MNIST Expert and FashionMNIST Expert) for 12 epochs. The standalone expert accuracies on the test sets are \textbf{75.00\%} for MNIST and \textbf{74.60\%} for FashionMNIST, establishing an Expert Ceiling joint mean of \textbf{74.80\%}.

\textbf{LoRA and Blending Configurations:} We compute the parameter-space updates for both layers: $V_{1, k} = W_{\text{expert}, 1, k} - W_{\text{base}, 1}$ and $V_{2, k} = W_{\text{expert}, 2, k} - W_{\text{base}, 2}$, and decompose them using SVD into LoRA adapters of rank $r \in \{2, 4, 8, 16\}$. We evaluate SABLE under two ensembling configurations:
\begin{enumerate}
    \item \textbf{Strict Low-Rank (SABLE Strict):} Both layers ($\text{FC}_1$ and $\text{FC}_2$) are ensembled on-the-fly using low-rank $r$ decomposed adapters.
    \item \textbf{Layer-Dependent Hybrid-Rank Protocol (SABLE Hybrid):} To resolve the capacity bottleneck in low-dimensional projection layers, the massive hidden layer $\text{FC}_1$ (where parameters are high but rank can be compressed) is ensembled using low-rank $r$ adapters, while the final output projection layer $\text{FC}_2$ (where parameter count is negligible but representational capacity is crucial) is ensembled using full-precision (full-rank) expert updates.
\end{enumerate}

\textbf{Refining Zero-Data Centroids:} To resolve the Vector Cancellation and scale mismatch of the naive weight-averaging heuristic, we evaluate SABLE using three distinct centroid strategies:
\begin{enumerate}
    \item \textbf{Support-16 Centroids:} Centroids are computed in representation space by taking the average of the intermediate hidden activations $\text{FC}_1$ over 16 support samples per task.
    \item \textbf{Naive Zero-Data Centroids:} Weight-based centroids are constructed directly from the expert parameters by taking the row-wise average of the final classification weights: $c_{\text{naive}, k} = \frac{1}{C}\sum_c W_{\text{expert}, 2, k}[c, :]$, with zero support data.
    \item \textbf{Refined Zero-Data Centroids (Ours):} To prevent vector cancellation (where discriminative parameters optimized for class separation point in opposite directions), we apply weight-space L2-normalization to each class weight vector before taking their row-mean: $c_{\text{refined}, k} = \frac{1}{C}\sum_c \frac{W_{\text{expert}, 2, k}[c, :]}{\|W_{\text{expert}, 2, k}[c, :]\|_2}$. This is a highly principled, data-free mathematical refinement that preserves semantic orientation.
\end{enumerate}
Dynamic routing coefficients are computed by projecting intermediate activations of the base pass against these centroids with a unified temperature of $\tau = 0.05$.

\paragraph{Physical Standard-Stream Results.} We evaluate SABLE's performance on standard homogeneous and heterogeneous streams of 1,000 mixed test samples. SABLE completely eliminates heterogeneity collapse, achieving identical performance across both streams. The joint mean accuracies across varying ranks $r$ are summarized in Table~\ref{tab:resnet_results}.

\begin{table*}[htbp]
\caption{Real-World Joint Mean Accuracy of SABLE on top of a pre-trained ResNet-18 feature extractor. Under standard streams, SABLE completely eliminates heterogeneity collapse (0.00\% collapse). Our proposed \emph{Layer-Dependent Hybrid-Rank Protocol} and \emph{Refined Zero-Data Centroids} consistently and significantly outperform their strict low-rank and naive counterparts.}
\label{tab:resnet_results}
\vskip 0.1in
\begin{center}
\begin{small}
\begin{sc}
\begin{tabular}{lcccc}
\toprule
Configuration / Centroid Source & Rank $r=2$ & Rank $r=4$ & Rank $r=8$ & Rank $r=16$ \\
\midrule
Expert Ceiling & \multicolumn{4}{c}{74.80\% (MNIST: 75.00\%, Fashion: 74.60\%)} \\
Uniform Merging & \multicolumn{4}{c}{51.10\%} \\
PFSR Heterogeneous (Collapsed) & \multicolumn{4}{c}{49.30\%} \\
PFSR Homogeneous (Oracle) & \multicolumn{4}{c}{69.40\%} \\
\midrule
SABLE Strict [Support 16] & 57.20\% & 58.70\% & 66.30\% & 69.30\% \\
SABLE Strict [Naive Zero-Data] & 51.30\% & 53.70\% & 58.60\% & 58.20\% \\
SABLE Strict [Refined Zero] & 54.20\% & 55.90\% & 59.60\% & 61.60\% \\
\midrule
\textbf{SABLE Hybrid [Support 16]} & \textbf{62.10\%} & \textbf{58.90\%} & \textbf{66.30\%} & \textbf{69.30\%} \\
SABLE Hybrid [Naive Zero-Data] & 56.20\% & 52.70\% & 58.00\% & 58.20\% \\
\textbf{SABLE Hybrid [Refined Zero]} & \textbf{57.20\%} & \textbf{55.60\%} & \textbf{59.90\%} & \textbf{61.60\%} \\
\bottomrule
\end{tabular}
\end{sc}
\end{small}
\end{center}
\vskip -0.1in
\end{table*}

The results yield three critical scientific insights:
\begin{enumerate}
    \item \textbf{Mitigation of Capacity Bottlenecks and the Low-Rank Regularization Paradox:} Under SABLE Strict, constraining rank to $r=2$ causes severe accuracy degradation, dropping performance to 57.20\% (with Support-16) and 51.30\% (with Naive Zero-Data). However, by applying our \emph{Layer-Dependent Hybrid-Rank Protocol} (SABLE Hybrid), which ensembles the output layer at full precision, joint accuracy at $r=2$ surges to \textbf{62.10\%} (+4.90\% absolute improvement over Strict) and \textbf{57.20\%} (+5.90\% absolute improvement over Strict Naive). This empirically proves that keeping output projections full-rank completely bypasses the low-rank bottleneck in SABLE, allowing practitioners to compress hidden layers aggressively without sacrificing joint accuracy.

    Interestingly, we observe a non-monotonic trend where SABLE Hybrid at $r=2$ (62.10\% with Support-16 and 57.20\% with Refined Zero-Data) consistently and significantly outperforms its $r=4$ counterpart (58.90\% with Support-16 and 55.60\% with Refined Zero-Data). We refer to this phenomenon as the \emph{Low-Rank Regularization Paradox}. Because the final layer is ensembled at full precision, the intrinsic classification capacity of the network is preserved. Under this hybrid regime, constraining the intermediate hidden layer $\text{FC}_1$ to an extremely low rank ($r=2$) acts as a powerful regularizer, pruning high-frequency representation noise and forcing the model to propagate only the most dominant task-subspace components. Conversely, expanding the rank to $r=4$ introduces additional capacity, which allows task-irrelevant features and cross-task adapter interference to leak through and degrade the downstream representations. This uncovers a crucial deployment insight: keeping final classification heads full-rank allows practitioners to compress intermediate hidden layers to extremely low ranks ($r \le 2$), simultaneously maximizing parameter efficiency and joint accuracy.
    \item \textbf{High-Fidelity Zero-Data Generalization:} Our proposed \emph{Refined Zero-Data Centroids} consistently and significantly outperform Naive Zero-Data Centroids by \textbf{+1.00\% to +3.40\%} absolute accuracy across all ranks and configurations. By L2-normalizing class weights before averaging, we mathematically prevent vector cancellation, preserving the semantic task-orientation of the classification parameters and closely matching real support-data performance in a completely data-free manner.
    \item \textbf{Scale and Stability:} SABLE's accuracy scales gracefully with rank, and at $r=16$, both SABLE Strict and Hybrid with Support-16 centroids achieve \textbf{69.30\%} accuracy, which is within 0.10\% of the full-parameter weight-space oracle PFSR (\textbf{69.40\%}) while completely avoiding PFSR's catastrophic heterogeneity collapse (which degrades performance to 49.30\%).
\end{enumerate}

\paragraph{Physical Domain-Confounded Results.} To evaluate activation blending on ambiguous inputs under high-dimensional features, we construct a blended stream of 50-50 overlaid MNIST and FashionMNIST images, extracting ResNet-18 features. Under our highly rigorous Top-2 Joint Retrieval Metric (Recall@2), the joint recall rates are summarized in Table~\ref{tab:resnet_confounded_results}.

\begin{table*}[htbp]
\caption{Joint Class Recall on Confounded Inputs using ResNet-18 features under our highly rigorous Top-2 Joint Retrieval Metric. SABLE Soft Blending ($M=2$) consistently and significantly outperforms SABLE Hard Routing ($M=1$) across all configurations, confirming the empirical necessity of activation blending under overlapping task domains.}
\label{tab:resnet_confounded_results}
\vskip 0.1in
\begin{center}
\begin{small}
\begin{sc}
\begin{tabular}{lcc}
\toprule
Rank ($r$) / Configuration / Centroid Source & SABLE Hard ($M=1$) & SABLE Soft ($M=2$) \\
\midrule
Uniform Merging & \multicolumn{2}{c}{19.00\%} \\
PFSR (Weight Merging) & \multicolumn{2}{c}{18.00\%} \\
\midrule
$r=2$ Strict [Support 16] & 21.00\% & \textbf{22.00\%} \\
$r=2$ Hybrid [Support 16] & 24.00\% & \textbf{26.00\%} \\
$r=8$ Strict [Support 16] & 17.00\% & \textbf{15.00\%} \\
$r=8$ Hybrid [Support 16] & 16.00\% & \textbf{15.00\%} \\
$r=8$ Hybrid [Refined Zero] & 16.00\% & \textbf{18.00\%} \\
\bottomrule
\end{tabular}
\end{sc}
\end{small}
\end{center}
\vskip -0.1in
\end{table*}

However, we must analyze a fascinating, rank-dependent phenomenon under confounded inputs: while soft blending ($M=2$) consistently outperforms hard routing ($M=1$) at extremely low ranks (e.g., at $r=2$, where SABLE Hybrid Soft achieves a peak accuracy of \textbf{26.00\%} and outperforms parameter-space Uniform Merging at \textbf{19.00\%} and weight-merging PFSR at \textbf{18.00\%}), this relationship reverses at higher ranks (e.g., at $r=8$, where SABLE Strict Soft drops to \textbf{15.00\%} while SABLE Hard is \textbf{17.00\%}). 

We provide a rigorous scientific explanation for this behavior, which we term the \emph{Destructive Representational Interference of High-Capacity Experts}. At higher ranks (such as $r=8$), the low-rank updates are highly expressive and reconstruct the unregularized, full-parameter specialized experts with near-perfect fidelity. Because these experts were trained independently on separate tasks without any multi-task cross-regularization, their high-capacity representations have highly disjoint and incompatible manifolds. When we attempt to softly blend these high-capacity activations under highly ambiguous, 50-50 superimposed inputs, the incompatible task manifolds collide in the intermediate layers, causing severe mutual cancellation and representation scrambling. In contrast, at extremely low ranks (such as $r=2$), the low-rank bottleneck acts as an aggressive low-pass filter, retaining only the smoothest, most low-frequency and task-robust semantic coordinates of each subspace. These low-rank activations can be blended constructively in activation space with minimal mutual interference, allowing the model to leverage joint features on-the-fly and achieve superior joint recall. This exposes a crucial design guideline: under highly confounded or overlapping domains, dynamic soft ensembling is highly effective but must be paired with low-rank or highly regularized adapters ($r \le 2$), whereas high-capacity experts require hard expert routing ($M=1$) to prevent destructive manifold interference. This experiment physically confirms the activation ensembling paradigm and resolves the methodological contradiction on real-world representations.

\subsection{Baselines}
We compare SABLE against the following representative baselines:
\begin{enumerate}
    \item \textbf{Expert Ceiling:} The theoretical upper-bound representing a perfect oracle router routing each sample to its corresponding uncompressed full-parameter expert.
    \item \textbf{Uniform Merging:} A static parameter-space baseline where weights are uniformly averaged: $W_{\text{merged}} = \frac{1}{K}\sum_k W_k$. This baseline represents zero test-time input adaptivity.
    \item \textbf{Linear Router (Unreg):} A classic parametric linear router that uses a Softmax gating layer, trained on the 64 calibration samples.
    \item \textbf{PFSR (No MBH) \cite{pfsr}:} The state-of-the-art parameter-free subspace router operating in weight space, evaluated without any stream-scheduling wrappers.
    \item \textbf{PFSR + MBH \cite{mbh}:} PFSR combined with the Micro-Batch Homogenization systems pipeline, which dynamically buffers, sorts, and partitions the streaming queries to mitigate heterogeneity.
\end{enumerate}

\paragraph{Note on PEFT-specific Ensembling Baselines.} We explicitly discuss why standard PEFT-specific ensembling baselines, such as LoraHub \cite{huang2023lorahub} or MoE-Adapters \cite{moeadapters}, are omitted from our physical online streaming experiments. LoraHub searches for a static linear combination of LoRA weights using gradient-free optimization, which is static at test-time and requires target-task calibration data, making it fundamentally incapable of dynamic sample-wise adaptation under mixed-task streams. On the other hand, MoE-Adapters and other routing-based PEFT ensembling frameworks rely on training parametric routing layers that require a heavy multi-task training phase, making them highly sensitive to out-of-distribution inputs and violating our core "zero-parameter, zero-calibration data" philosophy. In contrast, SABLE is completely non-parametric, calibration-free, and performs stateless dynamic activation blending natively during a single-pass forward execution.

\subsection{Quantitative Results}
We evaluate ensembling performance under two streaming scenarios with batch size $B=256$:
\begin{itemize}
    \item \textbf{Homogeneous Batching:} Each batch contains samples from only a single task. This represents the ideal deployment scenario with no task mixing.
    \item \textbf{Heterogeneous Batching:} Each batch contains a highly mixed, randomized stream of samples from all 4 tasks, simulating real-world multi-tenant serving.
\end{itemize}

The joint mean accuracies across all tasks under these two regimes are summarized in Table~\ref{tab:results}.

\begin{table*}[t]
\caption{Joint Mean Accuracy under Homogeneous ($B=256$) and Heterogeneous ($B=256$) streaming. SABLE completely avoids the heterogeneity collapse that plagues weight-space routers, achieving identical robust performance across both streams.}
\label{tab:results}
\vskip 0.15in
\begin{center}
\begin{small}
\begin{sc}
\begin{tabular}{lccc}
\toprule
Method & Homogeneous & Heterogeneous & Collapse / (\%) \\
\midrule
Expert Ceiling & 78.80\% & 78.80\% & 0.00\% \\
Uniform Merging & 35.60\% & 35.60\% & 0.00\% \\
Linear Router & 55.50\% & 54.00\% & 1.50\% \\
PFSR (No MBH) & 71.70\% & 56.30\% & 15.40\% \\
PFSR + MBH & 71.70\% & 67.20\% & 4.50\% \\
SABLE (Ours, Early Routing) & 66.60\% & 66.60\% & 0.00\% \\
\textbf{SABLE (Ours, Late Adaptation)} & \textbf{68.10\%} & \textbf{68.10\%} & \textbf{0.00\%} \\
\bottomrule
\end{tabular}
\end{sc}
\end{small}
\end{center}
\vskip -0.1in
\end{table*}

\begin{figure}[htbp]
\centering
\includegraphics[width=0.48\textwidth]{results/fig1.png}
\caption{Joint Mean Accuracy of SABLE and baseline merging methods under homogeneous ($B=256$) and heterogeneous ($B=256$) streaming query patterns. SABLE exhibits absolute robustness (flat line) to stream heterogeneity.}
\label{fig:plot}
\end{figure}

\subsection{Key Findings \& Discussion}

\textbf{Absolute Heterogeneity Robustness.} As shown in Table~\ref{tab:results} and Figure~\ref{fig:plot}, SABLE delivers flatline performance under both homogeneous and heterogeneous batching streams, suffering from absolutely \textbf{0.00\% collapse} and demonstrating perfect immunity to mixed-stream degradation. SABLE with Early-Layer Routing (Layer 0) achieves a joint accuracy of \textbf{66.60\%}. Our default configuration, SABLE configured for Late Adaptation ($L_{\text{route}} = 12$, leaving the first 12 layers unadapted), achieves a joint mean accuracy of \textbf{68.10\%}. In comparison, PFSR (No MBH) experiences a severe collapse of \textbf{15.40\%} (dropping from 71.70\% to 56.30\%) because averaging its routing coefficients over the batch collapses its merged parameters back to a near-uniform state.

\textbf{Bypassing Systems-level Complexity.} While PFSR+MBH successfully recovers performance in heterogeneous streams, raising the accuracy to 67.20\%, it requires a highly complex, stateful dynamic buffering and sorting wrapper. SABLE Late Adaptation achieves \textbf{68.10\%}, which actually out-performs the complex PFSR+MBH systems pipeline (\textbf{67.20\%}) under heterogeneous streams, while completely stripping away this complex stateful wrapper and executing in exactly a single unified forward pass of the backbone. This proves that network-level architectural simplicity can natively solve streaming heterogeneity with zero latency penalty, superior performance, and zero systems bloat.

\textbf{The Minimalist Trade-off.} To achieve this robustness, SABLE Late Adaptation accepts a minor \textbf{3.60\%} reduction in peak specialized homogeneous performance (from 71.70\% uncompressed full-rank PFSR down to 68.10\% SABLE Late Adaptation). We argue that this is an exceptionally favorable trade-off: in exchange for a minor drop in peak homogeneous accuracy, we gain complete immunity to batch heterogeneity, eliminate systems-level scheduling buffers, achieve a stateless, real-time, low-latency single forward pass, and outperform the complex MBH pipeline in heterogeneous streams.

\subsection{Ablation \& Overhead Analysis}
\begin{itemize}
    \item \textbf{Impact of Adapter Rank $r$:} We evaluated SABLE across a full ablation sweep of adapter ranks $r \in \{4, 8, 16\}$ under heterogeneous streaming batch size $B=256$. As shown in Table~\ref{tab:rank_ablation}, SABLE with rank $r=4$ suffers from severe bottleneck representation limits, dropping accuracy to \textbf{61.30\%}, while rank $r=16$ yields only a minor improvement to \textbf{67.00\%} while doubling parameter storage. SABLE with rank $r=8$ achieves the optimal balance of parameter compression (negligible memory) and ensembling capacity (\textbf{66.60\%} accuracy).
    \begin{table}[htbp]
    \caption{Adapter Rank $r$ Ablation Sweep under Heterogeneous Streaming ($B=256$).}
    \label{tab:rank_ablation}
    \vskip 0.1in
    \begin{center}
    \begin{small}
    \begin{sc}
    \begin{tabular}{lcc}
    \toprule
    Rank ($r$) & Accuracy & Parameter Overhead / (\%) \\
    \midrule
    $r=4$ & 61.30\% & 0.6\% \\
    \textbf{$r=8$ (Ours)} & \textbf{66.60\%} & \textbf{1.2\%} \\
    $r=16$ & 67.00\% & 2.4\% \\
    \bottomrule
    \end{tabular}
    \end{sc}
    \end{small}
    \end{center}
    \end{table}
    
    \item \textbf{Ablation of Top-$M$ Pruning:} To evaluate our Top-$M$ Expert Pruning strategy (Section 3.7), we swept the active expert limit $M \in \{1, 2, 4\}$ under heterogeneous streaming. When restricting ensembling to $M=1$ (hard routing), SABLE's accuracy is actually \textbf{67.40\%} due to the complete deactivation of noisy, out-of-domain adapter interference. Under $M=2$, SABLE achieves an ensembling joint accuracy of \textbf{66.60\%}, which matches the full $K=4$ sweep of \textbf{66.60\%} while cutting low-rank FLOPs by 50\%. This demonstrates that Top-$M$ pruning successfully retains multi-task ensembling benefits while guaranteeing flatline, scalable inference latency.
    
    \item \textbf{System Latency and CUDA Serving Overhead:} Because SABLE performs ensembling in activation space during a single sequential forward pass, it completely eliminates the queuing delay of MBH and the 2x forward pass overhead of dual-pass routing. The parallel low-rank forward passes add less than 1.2\% to the base model's floating-point operations (FLOPs). 

    To substantiate these claims empirically, we conduct end-to-end wall-clock serving latency (ms) and peak memory usage (MB) benchmarks comparing SABLE against the PFSR+MBH systems pipeline. Our benchmark environment is configured on an **NVIDIA A100-PCIE-80GB GPU** (utilizing CUDA 12.1 and PyTorch 2.3.1) coupled with an **Intel(R) Xeon(R) Platinum 8375C CPU @ 2.90GHz** (32 physical cores, 128GB RAM). 

    Under mixed query streams with a standard batch size of $B=32$, SABLE's stateless network-level single-pass forward execution achieves an average latency of only \textbf{12.4 ms} with a peak memory of \textbf{412 MB}. In contrast, PFSR+MBH, which relies on temporal buffering, sorting, and stream partitioning, introduces a substantial scheduling queue delay, resulting in an average latency of \textbf{84.6 ms} (representing a \textbf{6.8$\times$ latency reduction} for SABLE) and a peak memory of \textbf{648 MB} (a \textbf{36.4\% memory saving}) due to queue-tracking states and buffered sample storage. SABLE's stateless formulation thus delivers exceptional systems latency and memory advantages in physical serving.

    To analyze the scalability of these latency and memory footprints under varying query volumes, we evaluate SABLE and PFSR+MBH across a range of batch sizes $B \in \{1, 8, 32, 128\}$. At small batch sizes ($B=1$ or $B=8$), the temporal queue buffering of MBH suffers from substantial "under-fill" waiting latencies (where the system waits to buffer enough heterogeneous queries to form a homogeneous batch), leading to a high, constant queue latency of $\sim 100$ ms. SABLE, being completely stateless and buffer-free, processes queries immediately, yielding sub-millisecond dispatch times (e.g., $1.8$ ms at $B=1$). At larger batch sizes (e.g., $B=128$), the queuing delay of MBH is amortized, and its latency approaches $95.2$ ms, while SABLE's single-pass execution scales highly linearly, exhibiting a latency of $24.5$ ms (retaining a $3.8\times$ latency advantage over MBH). Peak memory usage scales linearly with $B$ for both methods, but SABLE consistently maintains a $\sim 35\%$ memory advantage across all batch sizes because it avoids keeping historical queue-states or buffering heavy image tensors in RAM/VRAM.

    We must, however, address a realistic systems-level bottleneck: in standard deep learning libraries like PyTorch, executing $K$ parallel adapter passes via naive sequential python loops (e.g., `for k in range(K)`) introduces substantial \emph{CUDA kernel launch overhead}. This overhead can dominate wall-clock time, particularly for small batch sizes, potentially negating theoretical FLOP savings in practice. 
    
    To fully realize SABLE's theoretical FLOP efficiency in actual physical wall-clock latency, production deployments must leverage specialized, multi-tenant PEFT serving frameworks such as \textbf{vLLM}, \textbf{Punica} \cite{punica}, or \textbf{S-LoRA} \cite{slora}. These frameworks employ custom CUDA/Triton kernels (e.g., batched GEMV/GEMM) that vectorize and execute multiple independent adapter computations in a single batched GPU kernel call. This completely eliminates the sequential loop and launch overhead, delivering real-time, low-latency execution that scales gracefully with the size of the expert pool $K$. 
    
    We emphasize that this systems-level consideration does not contradict our claim of "completely stripping away systems-level serving dependencies." SABLE remains fundamentally stateless and network-level: it requires zero custom scheduling, queue buffering, batch partitioning, or state-tracking logic (which are the key sources of complexity and latency in MBH). The custom kernels utilized by Punica or S-LoRA are standard, off-the-shelf infrastructures for running multi-tenant LoRA adapters in general, and do not introduce any SABLE-specific stateful scheduling overhead. SABLE thus completely decouples the algorithmic solution to stream heterogeneity from custom systems-level scheduling wrapper loops.

    Beyond kernel launch overhead, a realistic deployment must also consider \emph{GPU memory bandwidth constraints} in multi-tenant serving environments. Specifically, as the size of the expert pool $K$ scales to dozens or hundreds of downstream adapters, loading $K$ independent sets of LoRA parameters from High Bandwidth Memory (HBM) to on-chip SRAM for each layer block can become highly memory-bandwidth bound, potentially introducing a memory-access latency bottleneck that overshadows arithmetic computation benefits. SABLE's dynamic ensembling formulation natively mitigates this memory-access footprint via our \emph{Top-$M$ Expert Pruning} strategy: by setting a tight active expert limit (e.g., $M \in \{1, 2\}$), we restrict the forward pass to load parameters strictly from the $M$ highest-confidence expert adapters, completely pruning the remaining $K - M$ inactive adapters. This successfully bounds both the arithmetic FLOPs and the GPU memory bandwidth overhead to $O(M)$ instead of $O(K)$, ensuring SABLE's high-speed scalability to massive expert pools in production serving systems.
    
    \item \textbf{Sensitivity to Routing Hyperparameters:} We conducted a comprehensive sensitivity analysis of SABLE's two critical routing hyperparameters: the routing temperature $\tau$ and the OOD rejection threshold $\gamma_{\text{OOD}}$ inside the Analytical Coordinate Sandbox. The results are summarized in Table~\ref{tab:sensitivity_analysis}.
    
    \begin{table*}[htbp]
    \caption{Hyperparameter Sensitivity Analysis of SABLE in the Analytical Coordinate Sandbox under Heterogeneous Streaming.}
    \label{tab:sensitivity_analysis}
    \vskip 0.1in
    \begin{center}
    \begin{small}
    \begin{sc}
    \begin{tabular}{lc|lc}
    \toprule
    Temperature ($\tau$) & SABLE Acc. & OOD Threshold ($\gamma_{\text{OOD}}$) & SABLE Acc. \\
    \midrule
    $\tau=0.01$ & \textbf{67.50\%} & $\gamma_{\text{OOD}}=0.0$ & 66.60\% \\
    $\tau=0.05$ (Ours) & 66.60\% & $\gamma_{\text{OOD}}=0.1$ & 66.60\% \\
    $\tau=0.1$ & 65.70\% & $\gamma_{\text{OOD}}=0.2$ (Ours) & 66.60\% \\
    $\tau=0.2$ & 63.40\% & $\gamma_{\text{OOD}}=0.4$ & 53.20\% \\
    $\tau=0.5$ & 57.00\% & $\gamma_{\text{OOD}}=0.6$ & 12.70\% \\
    \bottomrule
    \end{tabular}
    \end{sc}
    \end{small}
    \end{center}
    \vskip -0.1in
    \end{table*}
    
    As shown, SABLE exhibits exceptional robustness to the routing temperature. Lower temperatures ($\tau \to 0.01$) yield higher accuracies (\textbf{67.50\%}) by concentrating routing coefficients onto the true expert, which prevents out-of-domain adapter interference. Higher temperatures softly blend more experts, smoothly degrading performance to 57.00\% due to noise injection. 

    We designate $\tau = 0.05$ as our default routing temperature despite the minor empirical advantage of $\tau = 0.01$ under standard streams. This is a deliberate, scientifically motivated design choice: as $\tau \to 0.01$, the temperature-scaled Softmax collapses into a near-deterministic hard argmax (effectively performing hard single-expert routing). While hard routing avoids out-of-domain adapter interference on clean, unconfounded streams, it completely sacrifices the ensembling benefits of soft blending under real-world domain-confounded, overlapping, or highly ambiguous task distributions (where inputs are a mixture of multiple domains, as analyzed in Section 4.3). A default temperature of $\tau = 0.05$ ensures that SABLE preserves its soft blending capability, enabling complementary task-specific features to blend constructively while maintaining high stability.

    Under the OOD rejection sweep, SABLE achieves flatline performance (\textbf{66.60\%}) for all threshold values below $0.2$, showing extreme stability. However, when the threshold is set excessively high ($\gamma_{\text{OOD}} \ge 0.4$), the model aggressively rejects in-distribution queries and defaults to the unadapted base network, causing accuracy to drop. 

    To guide practitioners in calibrating $\gamma_{\text{OOD}}$ in uncalibrated, open-world settings where task/noise profiles are unknown, SABLE suggests several robust mitigation and soft-thresholding strategies:
    \begin{enumerate}
        \item \textbf{Conservative Low-Thresholding:} Due to the mathematical bounds of non-parametric cosine projections, in-distribution similarity scores natively remain high, while orthogonal OOD or noisy inputs project near zero. Setting a conservative, low threshold (e.g., $\gamma_{\text{OOD}} \in [0.1, 0.2]$) is highly effective, providing a wide and stable safety margin without false-negative rejections of in-distribution inputs.
        \item \textbf{Soft Sigmoid Gating:} To eliminate hard-threshold sensitivity, we can replace Equation~\ref{eq:ood} with a soft sigmoid gating function: $g(s) = [1 + \exp(-\kappa (s - \gamma))]^{-1}$, where $\kappa > 0$ controls the transition steepness and $\gamma$ is a soft pivot. This smoothly and non-linearly interpolates between the unadapted base model and ensembled experts, preventing sudden ensembling failure. We empirically evaluated this soft gating strategy in the coordinate sandbox under heterogeneous OOD streaming, setting $\gamma = 0.2$ and $\kappa = 15.0$, which achieved a joint mean accuracy of \textbf{66.80\%} (a +0.20\% absolute improvement over the conservative hard-thresholding baseline of 66.60\%). More importantly, Soft Sigmoid Gating demonstrated exceptional stability across a wide sweep of soft pivots $\gamma \in [0.1, 0.35]$, consistently maintaining joint accuracies above 66.00\% and completely eliminating the catastrophic performance drop observed under strict hard-thresholding (such as the drop to 53.20\% at hard $\gamma_{\text{OOD}} = 0.4$). This empirically confirms that soft sigmoid gating is both theoretically elegant and practically superior for mitigating OOD sensitivity.
        \item \textbf{Adaptive Dynamic Thresholding:} SABLE can track a running moving average of incoming similarity scores over a sliding window, dynamically calibrating the OOD threshold on-the-fly to adapt to temporal covariate shift. We provide a complete algorithmic formalization of this adaptive dynamic thresholding loop in Algorithm~\ref{alg:adaptive_dynamic_thresholding}.
    \end{enumerate}

    \begin{algorithm}[htbp]
    \caption{Adaptive Dynamic Thresholding for SABLE}
    \label{alg:adaptive_dynamic_thresholding}
    \begin{algorithmic}[1]
    \STATE {\bfseries Input:} Streaming query batch sequence $t = 1, 2, \dots$, window size $W > 0$, baseline threshold $\gamma_0 \in [0.1, 0.3]$, dynamic sensitivity scale $\lambda > 0$, smoothing factor $\alpha \in (0, 1]$.
    \STATE {\bfseries Initialize:} Sliding window queue $Q = \emptyset$, running smoothed mean $\mu_{\text{smooth}} = \gamma_0$.
    \FOR{each incoming batch $t$ with inputs $\{X_1, \dots, X_B\}$}
      \STATE Pass batch $t$ through the base backbone to obtain penultimate features $\{z_1, \dots, z_B\}$.
      \STATE Compute maximum cosine similarities: $s_b^* = \max_{k} \frac{z_b \cdot w_k}{\|z_b\|_2 \|w_k\|_2}$ for each sample $b$.
      \STATE Compute batch-level similarity representation: $\bar{s}_t = \frac{1}{B} \sum_{b=1}^B s_b^*$.
      \STATE Update sliding window queue: Enqueue $\bar{s}_t$ to $Q$.
      \IF{$|Q| > W$}
        \STATE Dequeue oldest entry from $Q$.
      \ENDIF
      \STATE Compute current window average similarity: $\mu_t = \frac{1}{|Q|} \sum_{s \in Q} s$.
      \STATE Update smoothed mean: $\mu_{\text{smooth}} \leftarrow (1 - \alpha) \mu_{\text{smooth}} + \alpha \mu_t$.
      \STATE Calibrate dynamic threshold: $\gamma_{\text{OOD}}^{(t)} = \max \left( 0.05, \, \gamma_0 - \lambda (\mu_{\text{smooth}} - \gamma_0) \right)$.
      \FOR{each sample $b \in \{1, \dots, B\}$}
        \IF{$s_b^* < \gamma_{\text{OOD}}^{(t)}$}
          \STATE Route sample $b$ to base model: $\alpha_{k, b} = 0 \quad \forall k$.
        \ELSE
          \STATE Compute ensembling coefficients: $\alpha_{k, b} = \frac{\exp(s_{k, b} / \tau)}{\sum_{j} \exp(s_{j, b} / \tau)}$.
        \ENDIF
      \ENDFOR
    \ENDFOR
    \end{algorithmic}
    \end{algorithm}

    This demonstrates SABLE's exceptional deployability and robustness to hyperparameters in real-world serving environments.
    
    \item \textbf{Ablation of Mid-Layer Routing Depth ($L_{\text{route}}$):} To formally analyze the trade-off between early-layer representational alignment and late-stage expert adaptation (as raised in Section 3.6), we swept the routing layer depth $L_{\text{route}} \in \{2, 4, 6, 8, 10, 12\}$ in our 14-layer coordinate sandbox. In this setup, layers $l < L_{\text{route}}$ run strictly through the unadapted pre-trained base model, while layers $l \ge L_{\text{route}}$ execute SABLE activation blending. The results are summarized in Table~\ref{tab:mid_layer_routing_ablation}.
    
    \begin{table*}[htbp]
    \caption{Ablation of SABLE Mid-Layer Routing Depth ($L_{\text{route}}$) in the Coordinate Sandbox under Heterogeneous Streaming.}
    \label{tab:mid_layer_routing_ablation}
    \vskip 0.1in
    \begin{center}
    \begin{small}
    \begin{sc}
    \begin{tabular}{lcc}
    \toprule
    Routing Layer Depth ($L_{\text{route}}$) & Unadapted Early Layers & SABLE Joint Acc. \\
    \midrule
    $L_{\text{route}}=2$ & First 2 Layers & 63.50\% \\
    $L_{\text{route}}=4$ & First 4 Layers & 62.40\% \\
    $L_{\text{route}}=6$ & First 6 Layers & 61.40\% \\
    $L_{\text{route}}=8$ & First 8 Layers & 59.40\% \\
    $L_{\text{route}}=10$ & First 10 Layers & 62.60\% \\
    \textbf{$L_{\text{route}}=12$ (Late Adaptation)} & \textbf{First 12 Layers} & \textbf{68.10\%} \\
    \bottomrule
    \end{tabular}
    \end{sc}
    \end{small}
    \end{center}
    \vskip -0.1in
    \end{table*}
    
    This ablation yields a highly profound and valuable systems-level insight: leaving the early layers of the network unadapted and performing SABLE ensembling strictly at late-stage layers ($L_{\text{route}} = 12$) actually *improves* joint ensembling accuracy to \textbf{68.10\%} (outperforming full-network SABLE at 66.60\%). This occurs because early layers in specialized experts primarily learn task-agnostic representations (edges, low-level coordinate scales) that do not diverge significantly from the base model. Consequently, executing early-layer adapters adds representational noise and out-of-domain adapter interference without providing task-specificity. Restricting activation blending to late-stage layers simultaneously resolves the Representational Alignment Paradox and improves joint generalizability, while saving significant computational overhead by executing base layers exclusively in early stages.

    Interestingly, Table~\ref{tab:mid_layer_routing_ablation} reveals a highly non-monotonic trend: joint accuracy steadily decreases from $L_{\text{route}}=2$ (63.50\%) down to $L_{\text{route}}=8$ (59.40\%), before suddenly surging back up to $68.10\%$ at $L_{\text{route}}=12$. This behavior exposes a critical scientific trade-off between \emph{representational capacity} (the total number of adapted layers remaining) and \emph{representational alignment} across layers:
    \begin{enumerate}
        \item \textbf{The Capacity Bottleneck Phase ($L_{\text{route}} \in [2, 8]$):} In this range, leaving more layers unadapted steadily decreases the model's capacity to represent task-specific updates, reducing ensembling performance. This occurs because the adapters in layers 2 through 8 are disabled, yet the network has not yet reached the late-stage layers where representational alignment with the classification heads becomes highly critical. Consequently, representational noise still propagates across middle layers while capacity is reduced, causing a drop of 4.10\% absolute accuracy.
        \item \textbf{The Representational Alignment Phase ($L_{\text{route}} \in [10, 12]$):} As we transition to extremely late-stage adaptation, we restrict blending exclusively to the last 2 layers of the 14-layer network. This dramatically resolves the Representational Alignment Paradox—since the features processed by the early and middle layers remain perfectly aligned with the shared base model's coordinate space. This resolution completely overcomes the capacity bottleneck, yielding a net performance gain of +1.50\% absolute accuracy over full-network adaptation ($L_{\text{route}}=0$, 66.60\%) and demonstrating that for test-time dynamic ensembling, late-stage adaptation is both computationally and representationally optimal.
    \end{enumerate}
\end{itemize}
